\documentclass[aspectratio=169]{beamer}
\usetheme{Madrid}
\usecolortheme{whale}

\usepackage[utf8]{inputenc}
\usepackage[spanish]{babel}
\usepackage{amsmath,amssymb}
\usepackage{graphicx}
\usepackage{listings}
\usepackage{booktabs}

\lstset{
    language=Python,
    basicstyle=\ttfamily\footnotesize,
    keywordstyle=\color{blue}\bfseries,
    stringstyle=\color{red},
    commentstyle=\color{gray}\itshape,
    breaklines=true,
    showstringspaces=false
}

\title[INT210 - Sesión 0.7]{Sesión 0.7: Sets y Álgebra de Conjuntos en Alta Velocidad}
\subtitle{Módulo 0.2: Estructuras de Datos Nativas en Python}
\author[Diplomado en Ciencia de Datos]{Cátedra de Inteligencia Artificial y Ciencia de Datos}
\institute[INT210]{Machine Learning From Scratch}
\date{Semana 2}

\begin{document}

\frame{\titlepage}

\begin{frame}{Contenido de la Sesión}
    \tableofcontents
\end{frame}

\section{El Álbum de Figuritas y Conjuntos}

\begin{frame}{La Analogía: El Álbum de Figuritas sin Duplicados}
    \begin{columns}
        \column{0.55\textwidth}
        \begin{itemize}
            \item \textbf{Unicidad Automática:} No existen estampas duplicadas en un conjunto.
            \item \textbf{Membresía Instantánea:} Saber si una estampa existe toma tiempo $\mathcal{O}(1)$.
            \item \textbf{Intersección ($A \cap B$):} Figuritas repetidas que ambos amigos poseen para intercambiar.
        \end{itemize}
        \column{0.45\textwidth}
        \begin{alertblock}{Propiedad Hash}
            Los elementos de un \texttt{set} deben ser estrictamente inmutables (\textit{hashables}: \texttt{int}, \texttt{str}, \texttt{tuple}).
        \end{alertblock}
    \end{columns}
\end{frame}

\section{Operaciones Formales del Álgebra de Conjuntos}

\begin{frame}{Operaciones Fundamentales (Joel Grus, Cap. 2)}
    Sean dos conjuntos $A$ y $B$:
    \begin{center}
    \begin{tabular}{lll}
        \toprule
        \textbf{Operación Matemática} & \textbf{Sintaxis Python} & \textbf{Complejidad} \\
        \midrule
        \textbf{Unión ($A \cup B$)} & \texttt{A | B} & $\mathcal{O}(|A| + |B|)$ \\
        \textbf{Intersección ($A \cap B$)} & \texttt{A \& B} & $\mathcal{O}(\min(|A|, |B|))$ \\
        \textbf{Diferencia ($A \setminus B$)} & \texttt{A - B} & $\mathcal{O}(|A|)$ \\
        \textbf{Diferencia Simétrica ($A \Delta B$)} & \texttt{A \^{} B} & $\mathcal{O}(|A| + |B|)$ \\
        \textbf{Pertenencia ($x \in A$)} & \texttt{x in A} & $\mathcal{O}(1)$ \\
        \bottomrule
    \end{tabular}
    \end{center}
\end{frame}

\section{Aplicaciones en Ciberseguridad y Arte}

\begin{frame}[fragile]{Filtrado de Botnets y Deduplicación Cromática}
    \textbf{Filtrado Perimetral contra Lista Negra:}
\begin{lstlisting}
blacklist_total = feed_cert | feed_soc | feed_cloud
if ip_entrante in blacklist_total:  # Consulta en O(1)
    bloquear_paquete()
\end{lstlisting}
    \vspace{0.2cm}
    \textbf{Deduplicación de Paleta en Arte Digital:}
\begin{lstlisting}
paleta_unica = set(lista_pixeles_rgb)  # O(N) deduplicacion
\end{lstlisting}
\end{frame}

\begin{frame}{Conclusiones}
    \begin{itemize}
        \item El uso de \texttt{set} para comprobación de membresía es miles de veces más veloz que una búsqueda en \texttt{list}.
        \item El álgebra de conjuntos permite fusionar fuentes de inteligencia de amenazas con precisión matemática.
        \item \textbf{Siguiente Sesión (0.8):} Ingesta y Parseo de Archivos Delimitados "From Scratch".
    \end{itemize}
\end{frame}

\end{document}
